% !TEX root = ../AImath.v5.tex

\chapter{ 前言}

\section{从"困难重重"到"显而易见"：智能发展的历史辩证法}

在人工智能的发展历程中，我们见证了一个有趣的现象：昨天还被认为是"困难重重"的技术挑战，今天却变得"显而易见"。从早期的专家系统到今天的大语言模型，从手工特征工程到端到端的深度学习，每一次技术突破都遵循着这样的规律。

这种转变并非偶然，而是智能发展的内在逻辑。当我们真正理解了问题的本质，找到了合适的数学工具和计算方法，原本看似不可能的任务就会变得自然而然。正如牛顿所说："如果我看得更远，那是因为我站在巨人的肩膀上。"每一代人工智能研究者都在前人的基础上，将"困难"转化为"显而易见"。

\section{数学史：智能追求的发展主线}

数学不仅是描述世界的语言，更是智能发展的主线。从古希腊的几何学到现代的概率论，从牛顿的微积分到香农的信息论，每一个数学概念的诞生都为智能的理解和实现提供了新的工具。

线性代数让我们能够处理高维数据，微积分使得梯度下降成为可能，概率论为不确定性建模奠定了基础，而信息论则揭示了学习和压缩的本质联系。当我们回顾人工智能的发展历程时，会发现每一次重大突破都伴随着数学工具的创新或重新发现。

深度学习的成功，本质上是将这些数学工具有机结合的结果。神经网络的反向传播算法体现了微积分的链式法则，卷积神经网络利用了信号处理中的卷积理论，而注意力机制则借鉴了概率分布的思想。

\section{智能的层次：从模仿到创造}

智能的发展呈现出明显的层次性。最初级的智能是模仿——通过大量的数据学习人类的行为模式。这就像孩子学说话，通过不断重复和纠错来掌握语言规律。

更高层次的智能是理解——不仅知道"是什么"，还要知道"为什么"。这需要建立因果关系，理解概念之间的内在联系。当前的大语言模型在某种程度上展现了这种能力，它们不仅能够生成流畅的文本，还能进行推理和解释。

最高层次的智能是创造——能够产生全新的想法和解决方案。这要求系统不仅要掌握现有知识，还要能够重新组合、抽象和泛化，产生前所未有的洞察。虽然我们距离真正的创造性智能还有很长的路要走，但每一次技术进步都在向这个目标靠近。

\section{本书的目标与挑战}

本书的核心目标是帮助读者建立对人工智能的深层理解。我们不满足于简单的技术介绍，而是要揭示技术背后的数学原理、历史脉络和哲学思考。

我们面临的挑战是如何在保持严谨性的同时，让内容变得易于理解。数学是精确的语言，但也可能成为理解的障碍。我们的策略是：
\begin{itemize}
\item \emph{从直觉出发}：每个概念都从直观的例子开始，然后逐步深入到数学细节
\item \emph{历史视角}：通过历史发展的脉络来理解概念的来龙去脉
\item \emph{实践结合}：理论与实际应用相结合，让抽象概念变得具体
\item \emph{层次递进}：从简单到复杂，从具体到抽象，循序渐进地构建知识体系
\end{itemize}

我们希望读者在阅读本书后，不仅能够掌握人工智能的技术细节，更能够理解其背后的思想精髓，具备独立思考和创新的能力。

\section{致读者}

亲爱的读者，你即将踏上一段充满挑战但也充满收获的旅程。人工智能是一个快速发展的领域，新的技术和概念层出不穷。但是，正如我们在前面提到的，技术会过时，思想永恒；工具会更新，智慧长存。

本书的价值不在于教会你最新的技术细节——那些在几年后可能就会过时——而在于帮你建立正确的思维框架和分析方法。当你真正理解了智能的本质，掌握了数学的工具，培养了批判性思维，你就能够适应任何新的技术变化。

学习是一个主动的过程。我们鼓励你：
\begin{itemize}
\item \emph{质疑一切}：不要盲目接受任何观点，包括本书中的观点
\item \emph{动手实践}：理论必须与实践相结合才能真正掌握
\item \emph{深入思考}：不满足于表面的理解，要追问"为什么"
\item \emph{保持好奇}：对未知保持开放的心态，勇于探索新的可能性
\end{itemize}

技术会过时，思想永恒；工具会更新，智慧长存。愿这本书成为你探索智能奥秘的指南针，在"困难重重"的学习路上点亮一盏明灯，最终让那些看似高深的理论变得"显而易见"。

这本书代表了我们的尝试——让深度学习变得平易近人，教会人们概念、背景和思维。

我们以读者熟悉的案例说明方法的边界与取舍，避免抽象堆砌。通过具体的例子和实际的应用场景，我们希望将复杂的理论转化为可理解、可操作的知识。

\section{目标受众}

本书面向学生（本科生或研究生）、工程师和研究人员，他们希望扎实掌握深度学习的实用技术。因为我们从头开始解释每个概念，所以不需要过往的深度学习或机器学习背景。全面解释深度学习的方法需要一些数学和编程，但我们只假设读者了解一些基础知识，包括线性代数、微积分、概率和非常基础的Python编程。此外，在附录中，我们提供了本书所涵盖的大多数数学知识的复习。

大多数时候，我们会优先考虑直觉和想法，而不是数学的严谨性。有许多很棒的书可以引导感兴趣的读者走得更远。Bela Bollobas的《线性分析》 ([Bollobás, 1999](https://zh-v2.d2l.ai/chapter\_references/zreferences.htmlid13)) 
对线性代数和函数分析进行了深入的研究。 ([Wasserman, 2013]
(https://zh-v2.d2l.ai/chapter\_references/zreferences.htmlid178))
是一本很好的统计学指南。如果读者以前没有使用过Python语言，那么可以仔细阅读这个[Python教程]
%(http://learnpython.org/)。

\newpage